
Indeed, if \(G\) is an abelian variety of dimension \(d\), it is
known (cf.\ S. Lang, \emph{Abelian Varieties}, Chapter IV, \S3,
\SGARefLink{sga3:xvii:section:6:site:da2ef89f51}{Theorem 6}) that the group \({}_nG(k)\), for \(n\) prime to \(p\),
is isomorphic to
\[
  (\mathbb Z/n\mathbb Z)^{2d},
\]
and hence to \((\boldsymbol\mu_n)^{2d}\).  Thus
(iii)\(\Rightarrow\)(ii), while
(ii)\(\Rightarrow\)(i)\(\Rightarrow\)(iii) are clear.

\SGASubsection{4.6}{The General Case}
\SGARefTarget{sga3:xvii:subsection:4:6}

If \(G\) and \(H\) are two algebraic groups defined over an
algebraically closed field \(k\), we shall denote by \(P(G,H)\)
the property
\[
  \text{``there is no algebraic subgroup of \(G\) isomorphic to
  \(H\).''}
\]
We then obtain the following characterizations of unipotent
groups.

\SGAStatement{\SGAInlineRefTarget{sga3:xvii:theorem:4-6-1}Theorem 4.6.1}{Let \(G\) be an algebraic group
defined over an algebraically closed field \(k\) of characteristic
\(p\).  Then:
\begin{enumerate}
\item[(i)] \SGARefTarget{sga3:xvii:theorem:4:6:1:item:i}
  If \(G\) is smooth, affine, and connected, then
  \[
    \begin{split}
    G\text{ is unipotent}
      &\Longleftrightarrow
        \text{there exists \(n>1\) such that \(P(G,\boldsymbol\mu_n)\)
        holds}\\
      &\Longleftrightarrow P(G,\mathbb G_m)\text{ holds}.
    \end{split}
  \]
\item[(ii)] \SGARefTarget{sga3:xvii:theorem:4:6:1:item:ii}
  If \(G\) is smooth and connected, then
  \[
    G\text{ is unipotent}
      \Longleftrightarrow
    \text{there exists \(n\) prime to \(p\) such that
      \(P(G,\boldsymbol\mu_n)\) holds}.
  \]
\item[(iii)] \SGARefTarget{sga3:xvii:theorem:4:6:1:item:iii}
  If \(G\) is smooth, then
  \[
    G\text{ is unipotent}
      \Longleftrightarrow
    P(G,\boldsymbol\mu_n)\text{ holds for every prime \(n\ne p\)}.
  \]
\item[(iv)] \SGARefTarget{sga3:xvii:theorem:4:6:1:item:iv}
  If \(G\) is affine and connected and \(p>0\), then
  \[
    G\text{ is unipotent}
      \Longleftrightarrow P(G,\boldsymbol\mu_p)\text{ holds}.
  \]
\item[(v)] \SGARefTarget{sga3:xvii:theorem:4:6:1:item:v}
  If \(G\) is connected and \(p>0\), then
  \[
    \begin{split}
    G\text{ is unipotent}
      \Longleftrightarrow{}&
      \text{there exists \(n\) prime to \(p\) such that}\\[-2pt]
      &P(G,\boldsymbol\mu_n)\text{ and }
        P(G,\boldsymbol\mu_p)\text{ both hold}.
    \end{split}
  \]
\item[(vi)] \SGARefTarget{sga3:xvii:theorem:4:6:1:item:vi}
  For an arbitrary algebraic group \(G\),
  \[
    G\text{ is unipotent}
      \Longleftrightarrow
    P(G,\boldsymbol\mu_n)\text{ holds for every prime \(n\)}.
  \]
\end{enumerate}}
\SGARefTarget{sga3:xvii:theorem:4:6:1}

\SGAProof{Proof}
Assertion (i) follows from \SGARefLink{sga3:xvii:proposition:4-1-1}{4.1.1}, and (iv) from \SGARefLink{sga3:xvii:proposition:4-3-1}{4.3.1}.  We shall
prove (vi); assertions (ii), (iii), and (v) are proved in the same
way and are left to the reader.

Let \(G\) be an algebraic group.  If \(G\) is unipotent, then
\(P(G,\boldsymbol\mu_n)\) holds for every \(n>1\) by 2.4(i).
Conversely, suppose that \(P(G,\boldsymbol\mu_n)\) holds for every
prime \(n\), and let us prove that \(G\) is unipotent.  Write
\(G^0\) for the identity component of \(G\).  If \(G^0\) is
smooth, a classical theorem of
Chevalley\begingroup\renewcommand{\thefootnote}{*}\footnote[0]{%
\emph{Séminaire Bourbaki}, no.\ 145, 1956/57.}\endgroup
shows that \(G^0\) is an extension
of an abelian variety \(A\) by a smooth connected affine group
\(L\).  If \(G\) is not smooth---which entails \(p>0\)
(Exposé VI\textsubscript{B}, \SGARefLink{cor:sga3-VIB-1.6.1}{1.6.1})---there is an integer \(n>0\)
such that
\[
  G''=G^0/{}_{F^n}(G)
\]
is smooth (Appendix II, \SGARefLink{sga3:xvii:lemma:3-1}{3.1}).  Thus \(G''\) is an extension of an
abelian variety \(A\) by a smooth connected linear group \(L''\).
Let \(L\) be the inverse image of \(L''\) in \(G^0\).  It is still
affine and connected, since \({}_{F^n}(G)\) is radicial.  In every
case, therefore, \(G\) has a composition series
\[
  0\subset L\subset G^0\subset G
\]
such that \(L\) is affine and connected, \(G^0/L=A\) is an
abelian variety, and \(G/G^0\) is an étale group.

If \(P(G,\boldsymbol\mu_n)\) holds, then a fortiori
\(P(L,\boldsymbol\mu_n)\) holds, and hence \(L\) is unipotent
(\SGARefLink{sga3:xvii:proposition:4-1-1}{4.1.1} and \SGARefLink{sga3:xvii:proposition:4-3-1}{4.3.1}).  If \(A\ne0\), there are a prime number \(n\)
and a subgroup of \(A\) isomorphic to \(\boldsymbol\mu_n\) (\SGARefLink{sga3:xvii:subsection:4:5}{4.5}).
By \SGARefLink{sga3:xvii:subsection:5:1}{5.1} below, this subgroup lifts to \(G\), contradicting
\(P(G,\boldsymbol\mu_n)\).  Hence \(A=0\).  Finally, if
\(G/G^0\) is not unipotent, there are an integer \(q\) and a
subgroup of \(G/G^0\) isomorphic to \(\boldsymbol\mu_q\) (\SGARefLink{sga3:xvii:proposition:4:4:1}{4.4.1}).
As above, it follows that \(P(G,\boldsymbol\mu_q)\) does not hold.
Thus \(G/G^0\) is unipotent, and consequently so is \(G\).

\SGASection{5}{Extensions of Groups of Multiplicative Type by
Unipotent Groups}
\SGARefTarget{sga3:xvii:section:5}

\SGASubsection{5.1}{Statement of the Theorem}
\SGARefTarget{sga3:xvii:subsection:5:1}

\SGAStatement{\SGAInlineRefTarget{sga3:xvii:definition:5-1-0}Definition 5.1.0}{Let \(k\) be a field and \(G\) an
algebraic \(k\)-group.  Following the terminology introduced by
Rosenlicht (\emph{Questions of rationality for solvable algebraic
groups over non perfect fields}, Annali di Matematica 61 (1963)),
we shall say that \(G\) is \emph{\(k\)-solvable} if it satisfies
the following equivalent conditions:
\begin{enumerate}
\item[(i)] \SGARefTarget{sga3:xvii:definition:5:1:0:item:i}
  \(G\) has a composition series whose successive quotients are
  isomorphic to \(\mathbb G_a\).
\item[(ii)] \SGARefTarget{sga3:xvii:definition:5:1:0:item:ii}
  \(G\) has a characteristic composition series \((G_i)\) such
  that the successive quotients \(G_i/G_{i+1}\) are commutative
  and themselves have composition series whose successive
  quotients are isomorphic to \(\mathbb G_a\).
\end{enumerate}}
\SGARefTarget{sga3:xvii:definition:5:1:0}

The implication \SGARefLink{sga3:xvii:definition:5:1:0:item:i}{(i)}\(\Rightarrow\)\SGARefLink{sga3:xvii:definition:5:1:0:item:ii}{(ii)} is proved in
\emph{loc.\ cit.}  In fact, for the composition series \((G_i)\)
one may take the series introduced in the proof of \SGARefLink{sga3:xvii:lemma:3:9:item:ii}{3.9(ii)}.

\SGAStatement{\SGAInlineRefTarget{sga3:xvii:theorem:5-1-1}Theorem 5.1.1}{Let \(k\) be a field, \(U\) a
unipotent algebraic \(k\)-group, \(H\) a \(k\)-group of
multiplicative type, and \(E\) an algebraic \(k\)-group which is
an extension of \(H\) by \(U\), so that there is an exact sequence
\[
  1\longrightarrow U\longrightarrow E\longrightarrow H
    \longrightarrow1.
\]
Then:
\begin{enumerate}
\item[(i)] \SGARefTarget{sga3:xvii:theorem:5:1:1:item:i}
  The extension \(E\) is trivial in each of the following cases:
  \begin{enumerate}
  \item[(a)] \SGARefTarget{sga3:xvii:theorem:5:1:1:item:i:a}
    \(k\) is algebraically closed;
  \item[(b)] \SGARefTarget{sga3:xvii:theorem:5:1:1:item:i:b}
    \(k\) is perfect and one of the groups \(U\) or \(H\) is
    connected;
  \item[(c)] \SGARefTarget{sga3:xvii:theorem:5:1:1:item:i:c}
    \(U\) is \(k\)-solvable;
  \item[(d)] \SGARefTarget{sga3:xvii:theorem:5:1:1:item:i:d}
    \(U\) is smooth and \(H\) is connected.
  \end{enumerate}
\item[(ii)] \SGARefTarget{sga3:xvii:theorem:5:1:1:item:ii}
  If \(H'\) and \(H''\) are two lifts of \(H\) in \(E\), then
  they are conjugate by an element of \(U(k)\) in each of the
  following cases:
  \begin{enumerate}
  \item[(a)] \SGARefTarget{sga3:xvii:theorem:5:1:1:item:ii:a}
    \(k\) is algebraically closed and \(H\) is smooth;
  \item[(b)] \SGARefTarget{sga3:xvii:theorem:5:1:1:item:ii:b}
    \(k\) is algebraically closed and \(U\) is smooth.
  \end{enumerate}
\end{enumerate}
(During the proof we shall indicate other cases in which the
conclusion of \SGARefLink{sga3:xvii:theorem:5:1:1:item:ii}{(ii)} is true without assuming that \(k\) is
algebraically closed.)}
\SGARefTarget{sga3:xvii:theorem:5:1:1}

If \(U\) is an algebraic group (respectively, a commutative
algebraic group) defined over a field \(k\), we denote by
\(H^1(k,U)\) (respectively, by \(H^i(k,U)\), \(i\geq0\)) the
pointed set (respectively, the \(i\)-th group) of Galois
cohomology of \(k\) with values in \(U\); see J.-P. Serre,
\emph{Cohomologie galoisienne}, Lecture Notes in Mathematics
no.\ 5, Springer.
